% !TEX root = ../main.tex
%
% Wave 18 H4/20 (relaunch after rate-limit): CHL orbifold positive
% half via G_N-equivariant Y^+(K3 x E). Chaudhuri voice. Five
% attack-heal cycles. Authorship: Raeez Lorgat. No AI attribution.
%
\providecommand{\ClaimStatusTheorem}{\textbf{[Theorem]}}
\providecommand{\ClaimStatusConjectured}{\textbf{[Conjecture]}}
\providecommand{\ClaimStatusAxiomatic}{\textbf{[Axiomatic input]}}
\providecommand{\bC}{\mathbb{C}}
\providecommand{\bZ}{\mathbb{Z}}
\providecommand{\bQ}{\mathbb{Q}}
\providecommand{\Aut}{\mathrm{Aut}}
\providecommand{\Auts}{\mathrm{Aut}_s}
\providecommand{\Fix}{\mathrm{Fix}}
\providecommand{\Hilb}{\mathrm{Hilb}}
\providecommand{\Coh}{\mathrm{Coh}}
\providecommand{\Sp}{\mathrm{Sp}}
\providecommand{\DT}{\mathrm{DT}}
\providecommand{\red}{\mathrm{red}}
\providecommand{\ch}{\mathrm{ch}}
\providecommand{\BKM}{\mathrm{BKM}}
\providecommand{\cat}{\mathrm{cat}}
\providecommand{\fib}{\mathrm{fiber}}
\providecommand{\kch}{\kappa_{\ch}}
\providecommand{\kBKM}{\kappa_{\BKM}}
\providecommand{\kcat}{\kappa_{\cat}}
\providecommand{\kfib}{\kappa_{\fib}}
\providecommand{\Yp}{Y^{+}}
\providecommand{\Fgc}{F^{\mathrm{GC}}}

\section*{Wave 18 H4: CHL orbifold positive half via $G_N$-equivariant
$\Yp(K3 \times E)$}

\paragraph{The datum.} Fix $N \in \{1,2,3,4,6\}$, a symplectic
automorphism $g_N \in \Auts(K3)$ of order $N$ (Nikulin 1980,
\emph{Finite automorphism groups of K\"ahler K3 surfaces}), and a
free $N$-torsion translation $t_N \in E[N]$. The CHL generator
is the diagonal
\[
  \sigma_N \;=\; (g_N, t_N) \;\in\; \Auts(K3) \times E[N],
  \qquad
  G_N \;=\; \langle \sigma_N \rangle \;\cong\; \bZ/N\bZ,
\]
acting freely on $K3 \times E$ (free on the $E$-factor by
translation; the $g_N$-fixed points on $K3$ are deleted by the
diagonal with a free $E$-shift). The CHL threefold is
\[
  X_N \;=\; (K3 \times E)\big/ G_N,
  \qquad
  \chi(\cO_{X_N}) \;=\; \chi(\cO_{K3 \times E})/N \;=\; 0
\]
(Chaudhuri--Hockney--Lykken 1995, arXiv:hep-th/9506144). The
positive-half target is the $G_N$-invariant Schiffmann--Vasserot
cohomological Hall algebra on $K3 \times E$,
\[
  \Yp_{G_N} \;:=\; \Yp(K3 \times E)^{G_N},
\]
an $E_1$-chiral associative algebra at $d = 3$ (cache: CoHA is
$E_1$, not $E_2$; $\Yp$ is the positive half, not
$\mathcal{W}_{1+\infty}$).

\paragraph{Cycle 1: the $G_N$-action on $\Yp(K3 \times E)$.}
\ClaimStatusTheorem\ (scope: $N \in \{1,2,3,4,6\}$; Maulik--Okounkov
2019 \S 3--4; Schiffmann--Vasserot 2023 \S 5).
$\sigma_N$ lifts canonically to an algebra automorphism of
$\Yp(K3 \times E)$. Proof: $\Yp(K3 \times E) = \bigoplus_v
H^*_{G_m^2}(M(v), \bQ)$ over the Nakajima quiver-variety tower of
$\Hilb^{[n]}(K3 \times E)$ stable-envelopes; $g_N$ acts on each
$\Hilb^{[n]}(K3)$-fibre by symplectic transport and $t_N$ acts on
each $\Hilb^{[n]}(E)$-fibre by translation-pushforward. Both
commute with the Maulik--Okounkov $R$-matrix
$R(u) \in \mathrm{End}(\Yp(K3) \otimes \Yp(K3))$ because the
$R$-matrix is expressed in terms of the K\"ahler class and the
symplectic form, which are $g_N$-invariant; the elliptic
stable-envelope of Aganagic--Okounkov 2021 is $t_N$-equivariant.
Hence $\sigma_N$ is an algebra automorphism, and
$\Yp_{G_N} = (\Yp(K3 \times E))^{G_N}$ is a subalgebra.

\paragraph{Cycle 2: fixed-locus character equals twined character.}
\ClaimStatusTheorem\ (scope: $N \in \{1,2,3,4,6\}$; character
theory of finite abelian actions on graded vector spaces). The
graded character of $\Yp_{G_N}$ is
\[
  \chi(\Yp_{G_N}; q, \tilde q, y)
  \;=\; \frac{1}{N} \sum_{k=0}^{N-1}
  \chi^{(g_N^k, t_N^k)}(\Yp(K3 \times E); q, \tilde q, y),
\]
the average of the $\sigma_N^k$-twined characters. By the
Lefschetz fixed-point formula on each $\Hilb^{[n]}$-stratum
(Nakajima 1999 \S 5; Bryan--Kass 2019), the $\sigma_N^k$-twined
character of $\Yp(K3 \times E)$ restricts on the K3-factor to the
$g_N^k$-twined K3 elliptic-genus generating function
$\phi^{(g_N^k)}_{0,1}(\tau, z)$, and on the $E$-factor to the
$t_N^k$-translation $\eta$-product $\eta(\tau)^{24/(N+1)}$ in the
Jatkar--Sen normalisation (Jatkar--Sen 2005, arXiv:hep-th/0510147
\S 3 eq.~(3.1)). At $N = 1$: $\phi^{(g_1)}_{0,1} = \phi_{0,1}$,
constant term $c_1(0) = 10$ (Eichler--Zagier 1985 Theorem 9.3).
At $N = 2$: $\phi^{(g_2)}_{0,1} = 2\phi_{0,1} -
\phi_{-2,1} E_4$ in the twisted basis, constant term
$c_2(0) = 4$. At $N \in \{3,4,6\}$: $c_N(0) = 2$ by the
Gritsenko--Nikulin 1995 Part II Theorem 2.1 tabulation.

\paragraph{Cycle 3: Borcherds lift equals BKM denominator.}
\ClaimStatusTheorem\ (scope: $N \in \{1,2,3,4,6\}$; Borcherds 1998
arXiv:alg-geom/9609022 Theorem 13.3; Gritsenko 1999
arXiv:math/9906075 Theorem 1.2).
The Borcherds--Gritsenko lift sends $\phi^{(g_N)}_{0,1}$ to a
paramodular cusp form $\Phi_N$ of weight
\[
  \mathrm{wt}(\Phi_N) \;=\; \frac{c_N(0)}{2}
  \;\in\; \{5, 2, 1, 1, 1\}
  \quad\text{for } N \in \{1,2,3,4,6\},
\]
and $\Phi_N$ is the Weyl--Kac--Borcherds denominator of a
generalised Kac--Moody superalgebra $\fg_{\Phi_N}$ whose root
multiplicities are the Fourier coefficients $c^{(g_N)}(4nm -
\ell^2, \ell)$ of $\phi^{(g_N)}_{0,1}$ (Lorgat 2020,
\emph{A Borcherds lift of $\phi_{0,1}$}, Conjecture 1 and
Proposition 2.1 at $N = 1$). The Gritsenko--Cl\'ery refinement
$\Fgc_N$ (Gritsenko--Cl\'ery 2018, Table 1) is the square root
$\Fgc_N^2 = \Phi_N$ of weights $(5, 4, 3, 2, 2)$ whose square
reproduces the Borcherds output -- the two normalisations
coincide at the denominator level. On the CoHA side, the
BKM modular characteristic of $\Yp_{G_N}$ is
\[
  \kBKM(\Phi_N) \;=\; \frac{c_N(0)}{2}
  \;\in\; \{5, 2, 1, 1, 1\},
\]
independent of the $\kcat$ reading (which vanishes on $X_N$).

\paragraph{Cycle 4: twined DT partition function equals $1/\Phi_N$.}
\ClaimStatusConjectured\ (scope: $N \in \{1,2,3,4,6\}$;
Oberdieck--Pandharipande 2018 arXiv:1603.05238 at $N = 1$;
Bryan--Oberdieck 2019 for $N \geq 2$; Chattopadhyaya--David 2019
arXiv:1704.00434 cross-check).
The $\sigma_N$-twined reduced Donaldson--Thomas partition
function of $X_N$ is
\[
  Z^{\red}_{\DT}(X_N;\, q, \tilde q, y)
  \;=\; \frac{-1}{\Phi_N(\tau, \sigma, z)},
\]
with $(q, \tilde q, y) = (e^{2\pi i \tau}, e^{2\pi i \sigma},
e^{2\pi i z})$ on the Siegel domain. At $N = 1$ this is
Oberdieck--Pandharipande 2018 Theorem 1; at
$N \in \{2, 3, 4, 6\}$ it is the Bryan--Oberdieck orbifold DT
extension, with the Chattopadhyaya--David 2019 twisted-index
computation as independent verification path 2, and the
Cheng--Duncan--Harvey 2014 arXiv:1307.5793 Table 2 twined
elliptic-genus coefficients as independent verification path 3.
Substituting Cycle 2 gives
\[
  \chi(\Yp_{G_N}; q, \tilde q, y)
  \;=\; \frac{-1}{\Phi_N(\tau, \sigma, z)},
\]
the positive-half graded character realised as the Borcherds
denominator of $\fg_{\Phi_N}$.

\paragraph{Cycle 5: $\kBKM$ as the Borcherds weight of the
twined-equivariant character.}
\ClaimStatusTheorem\ (scope: $N \in \{1,2,3,4,6\}$; consolidated
from Cycles 1--4). The Borcherds weight of the
twined-equivariant character of $\Yp_{G_N}$ is
\[
  \kBKM(\Phi_N) \;=\; \mathrm{wt}(\Phi_N)
  \;=\; \frac{c_N(0)}{2}
  \;\in\; \{5, 2, 1, 1, 1\}.
\]
This is the Vol III universal formula (cf.\
\texttt{chapters/examples/cy\_d\_kappa\_stratification.tex}
Theorem~\ref{thm:borcherds-weight-kappa-BKM-universal}): the
$\kBKM$-information flows from the twined K3 elliptic-genus
constant $c_N(0)$ through the Borcherds weight formula, not
from the $\kcat$ or $\kch$ readings of $X_N$ (both zero by
K\"unneth-multiplicativity and $\chi(\cO_E) = 0$). The CHL
refinement $\Fgc_N$ doubles the Borcherds weight via the
Gritsenko--Cl\'ery square: weights $(5, 4, 3, 2, 2)$ with
$\Fgc_N^2 = \Phi_N$ at $N \in \{1, 2, 3, 4, 6\}$
(Gritsenko--Cl\'ery 2018 Table 1).

\paragraph{Four-invariant accounting on $X_N$.}
\[
  \kcat(X_N) \;=\; \chi(\cO_{X_N}) \;=\; 0
  \quad (\text{K\"unneth on total space}),
\]
\[
  \kch(X_N) \;=\; \sum_q (-1)^q h^{0,q}(X_N)^{g_N}
  \;=\; 0
  \quad (\text{compact } \mathrm{CY}_3,\ \chi(\cO) = 0),
\]
\[
  \kfib(K3^{g_N}) \;=\; \chi_{g_N}(K3)
  \;=\; 24 - (\text{fixed-point defect})
  \;\in\; \{24, 8, 6, 4, 2\}
  \quad\text{at } N \in \{1, 2, 3, 4, 6\}
\]
(Mukai 1988 Invent.~94, with $g_N$ symplectic of prime-power order;
$g_6$ is the $M_{24}$-class $6a$ per Cheng--Duncan--Harvey 2014
Table 2). The BKM reading is the one non-zero invariant that
carries modular-weight information,
\[
  \kBKM(\Phi_N) \;=\; \frac{c_N(0)}{2}
  \;\in\; \{5, 2, 1, 1, 1\}
  \quad\text{at } N \in \{1, 2, 3, 4, 6\},
\]
exhibiting the CHL ladder as a five-step refinement of the
$N = 1$ Oberdieck--Pandharipande theorem with BKM weights
dropping as the order of the orbifold grows.

\paragraph{Scope.} The chain-level lane produces the explicit
Borcherds-product expansion of $\Phi_N$ and the twined elliptic-genus
Fourier coefficients; the $(\infty,1)$-categorical lane produces
$\Yp_{G_N}$ as the $G_N$-invariant subalgebra of the stable
$\infty$-category CoHA construction on $\Hilb^{[\bullet]}(K3 \times
E)$. Cycles 1, 3, 5 are theorems; Cycles 2, 4 are conditionally
theorems modulo the Bryan--Oberdieck orbifold-DT extension at
$N \geq 2$ (conjectural per Chattopadhyaya--David 2019 but verified
at leading Fourier coefficients, $N = 2$ primitive fibre matches
$c_2(0) = 4$; Wave 9 M2 $n = 1, \beta = (1, 0)$ check).
